% Generated draft: requires mathematical and source review.
\documentclass[11pt]{article}
\usepackage[margin=25mm]{geometry}
\usepackage{fontspec}
\setmainfont{Latin Modern Roman}
\usepackage{amsmath,amssymb,mathtools,unicode-math}
\setmathfont{Latin Modern Math}
\usepackage{enumitem,needspace,xurl,hyperref}
\setlength{\parindent}{0pt}
\setlength{\parskip}{5pt}
\setlength{\emergencystretch}{3em}
\newcommand{\N}{\mathbb{N}}
\newcommand{\Z}{\mathbb{Z}}
\newcommand{\Q}{\mathbb{Q}}
\newcommand{\R}{\mathbb{R}}
\newcommand{\C}{\mathbb{C}}
\newcommand{\F}{\mathbb{F}}
\newcommand{\A}{\mathbb{A}}
\newcommand{\PP}{\mathbb{P}}
\newcommand{\E}{\mathbb{E}}
\newcommand{\eps}{\varepsilon}
\newcommand{\dd}{\,\mathrm{d}}
\begin{document}
\section*{UnsolvedMath ID: 30005402}
\textbf{Source problem code:} catalogue entry
\subsection{Definitions and mathematical statement}
Let \(W_{\Q_p}\) denote the Weil group of \(\Q_p\). In the geometric setting intended by the source, write \(G\) for the reductive group, \(\widehat G\) for its Langlands dual group, \(\operatorname{Bun}_G\) for the stack of \(G\)-bundles, and \(D(\operatorname{Bun}_G)\) for the sheaf category under consideration. The source does not specify all hypotheses on \(G\), the coefficient conventions, or the precise definition of this category; the formulation below is conditional on supplying these data.

For each dominant cocharacter \(\mu\) in the indexing set intended by the source, let \(r_\mu\colon\widehat G\to\operatorname{GL}(V_\mu)\) be the associated highest-weight representation. The source describes Hecke operators
\[
T_\mu\colon D(\operatorname{Bun}_G)\longrightarrow D(\operatorname{Bun}_G)^{B W_{\Q_p}},
\]
where the target denotes objects equipped with a continuous \(W_{\Q_p}\)-action. For a semisimple parameter written in the source as
\[
\phi\colon W_{\Q_p}\longrightarrow\widehat G(\Q_\ell),
\]
a Hecke eigensheaf with eigenvalue \(\phi\) is an object \(\mathcal S_\phi\in D(\operatorname{Bun}_G)\) equipped with \(W_{\Q_p}\)-equivariant isomorphisms
\[
T_\mu(\mathcal S_\phi)\simeq
\mathcal S_\phi\otimes (r_\mu\circ\phi)
\qquad\text{for every such }\mu.
\]
Here the tensor factor on the right is the Weil representation obtained from \(r_\mu\circ\phi\). The required coherence conditions and normalizations of these isomorphisms are not specified in the supplied excerpt.

\textbf{Problem.} Extend the eigensheaf construction attributed to Fargues from supercuspidal parameters to general semisimple \(L\)-parameters \(\phi\): construct and describe the intended objects \(\mathcal S_\phi\) and their Hecke eigenisomorphisms, beginning with parameters induced from a torus. The source describes the supercuspidal case using parameters that do not factor through a proper Levi subgroup of \(\widehat G\). Its phrase ``induced from a torus'' requires clarification: a possible intended meaning is a toral parameter composed with an appropriate map to the \(L\)-group of \(G\), but no torus, map, or induction convention is specified.

For context, the source indexes Harder--Narasimhan strata \(j_b\colon\operatorname{Bun}_G^b\hookrightarrow\operatorname{Bun}_G\) by the Kottwitz set \(B(G)\), and relates sheaves on a stratum to smooth representations of the \(\sigma\)-centralizer \(J_b(\Q_p)\). It refers to extensions by zero of representations in packets \(\Pi_\phi(J_b)\) on basic strata in the supercuspidal construction, and to an implication for a Kottwitz conjecture. The relevant displayed expression and the surrounding implication are truncated, so neither is asserted here as a precise formula or theorem.

\textbf{Unresolved qualifications.} The eigenisomorphisms alone allow the zero object. The source does not supply the nontriviality, packet, support, or normalization requirements needed to identify the intended construction. It also leaves unresolved the full \(L\)-group convention for nonsplit groups and the precise meaning of semisimplicity. Consequently, this is a conditional formulation of the research objective, requiring human mathematical review.
\subsection{Short English statement}
Extend Fargues's geometric eigensheaf construction to general semisimple Langlands parameters, starting with parameters induced from a torus. For each intended parameter \(\phi\), seek a specified sheaf \(\mathcal S_\phi\) whose Hecke transforms have eigenvalues \(r_\mu\circ\phi\) for every relevant dominant cocharacter \(\mu\). The supplied excerpt leaves essential conventions and characterization requirements unresolved. Its dated curation reports partial results for sufficiently generic toral parameters, with compatibility and full parameter generality still open; this assessment has not been independently verified.
\subsection{Sources}
\begin{itemize}
\item Supplied local UnsolvedMath registry record, permanent ID \(30005402\), title ``Fargues Eigensheaves for General Langlands Parameters,'' exact source problem code \texttt{catalogue entry}. No individual proposer is identified.
\item Supplied report citation: \emph{Arithmetic of Shimura Varieties} (2023), Oberwolfach Reports, pp.~261--326, DOI \texttt{10.4171/owr/2023/5}. \url{https://doi.org/10.4171/owr/2023/5}. The record attributes the original vision to Fargues, with an unresolved internal reference [4], and mentions the Kottwitz conjecture as [8, Conjecture 7.3]; the corresponding bibliography entries and precise problem location are not supplied.
\item Supplied literature source: Hamann, ``geometric Eisenstein series.'' \url{https://arxiv.org/abs/2209.08175}.
\item Supplied literature source: Fargues--Scholze, ``geometrization.'' \url{https://arxiv.org/abs/2102.13459}.
\item Literature assessment dated exactly \texttt{2026-08-22}: the supplied curation labels the problem partially solved and reports partial answers through normalized geometric Eisenstein functors for sufficiently generic toral parameters, with compatibility and full parameter generality still open. This is supplied dated curation, not an independently verified literature assessment.
\item The registry labels the statement \texttt{corrected\_verified} and says it was reconstructed and checked against the official report. That verification claim is supplied metadata; no independent source consultation or mathematical verification is claimed here.
\end{itemize}
\end{document}
